\subsection*{Questão 02}

\begin{questionbox}{02} 
  Sobre os autovalores e autovetores da matriz $A$ abaixo, é correto afirmar que:

  \[
  A = \begin{pmatrix}
  4 & 1 \\
  0 & 2
  \end{pmatrix}
  \]
\end{questionbox}

\begin{solutionbox}
  O polinômio característico de $A$ é dado por
  \[
  \det(A - \lambda I) = 0.
  \]
  Assim,
  \[
  \det \begin{pmatrix}
  4 - \lambda & 1 \\
  0 & 2 - \lambda
  \end{pmatrix}
  = (4 - \lambda)(2 - \lambda) = 0.
  \]
\end{solutionbox}

\begin{solutionbox}
  Resolvendo a equação característica:
  \[
  (4 - \lambda)(2 - \lambda) = 0
  \;\Rightarrow\;
  \lambda_1 = 4, \quad \lambda_2 = 2.
  \]

  Portanto, a matriz possui dois autovalores reais e distintos.
\end{solutionbox}

\begin{solutionbox}
  Para $\lambda_1 = 4$, resolvemos $(A - 4I)\mathbf{v} = 0$:
  \[
  \begin{pmatrix}
  0 & 1 \\
  0 & -2
  \end{pmatrix}
  \begin{pmatrix}
  x \\ y
  \end{pmatrix}
  =
  \begin{pmatrix}
  0 \\ 0
  \end{pmatrix}
  \;\Rightarrow\; y = 0.
  \]
  Logo, um autovetor associado é
  \[
  \mathbf{v}_1 = \begin{pmatrix}1 \\ 0\end{pmatrix}.
  \]
\end{solutionbox}

\begin{solutionbox}
  Para $\lambda_2 = 2$, resolvemos $(A - 2I)\mathbf{v} = 0$:
  \[
  \begin{pmatrix}
  2 & 1 \\
  0 & 0
  \end{pmatrix}
  \begin{pmatrix}
  x \\ y
  \end{pmatrix}
  =
  \begin{pmatrix}
  0 \\ 0
  \end{pmatrix}
  \;\Rightarrow\; 2x + y = 0.
  \]
  Assim, $y = -2x$ e um autovetor associado é
  \[
  \mathbf{v}_2 = \begin{pmatrix}1 \\ -2\end{pmatrix}.
  \]
\end{solutionbox}